%% ============================================================
%%  §5  ProgramSpec: A Shared Specification Language
%% ============================================================

\section{ProgramSpec: A Shared Specification Language}
\label{sec:programspec}

A central claim of this work is that a single specification language can
describe properties of programs from different source languages. This section
defines that language — \ProgramSpec — and demonstrates that its 17 proved
combinators are sufficient to state and verify all correctness properties in
the library.

\subsection{Motivation}

Consider two implementations of \texttt{max(a, b)}:
\begin{enumerate}[leftmargin=1.5em]
  \item A Rust function compiled to an \texttt{LLBCFunDef} and run via
        \texttt{evalFun}.
  \item A TypeScript function compiled to a \texttt{TSFunDef} and run via
        \texttt{evalTSFun}.
\end{enumerate}

Both should satisfy the same specification: ``if $a \ge b$, the result is
$a$; otherwise the result is $b$.'' Writing this specification once and
applying it to both implementations requires a specification language that
is \emph{parametric in the execution mechanism}. \ProgramSpec achieves this
by specifying properties of \RustM computations, which both evaluators return.

\subsection{The ProgramSpec Inductive}

\begin{lstlisting}[caption={ProgramSpec inductive type.},label={lst:progspec}]
inductive ProgramSpec (s a : Type) where
  | pureOutput : (a -> Prop) -> ProgramSpec s a
  | nocrash    : ProgramSpec s a
  | both       : ProgramSpec s a -> ProgramSpec s a -> ProgramSpec s a
  | withFuel   : Nat -> ProgramSpec s a -> ProgramSpec s a
  | precond    : Prop -> ProgramSpec s a -> ProgramSpec s a
\end{lstlisting}

The constructors have the following meanings:

\begin{itemize}[leftmargin=1.5em]
  \item \texttt{pureOutput P}: the computation returns some value $v$ such
        that $P\,v$ holds and does not modify the state.
  \item \texttt{nocrash}: the computation does not return \texttt{.error}.
  \item \texttt{both P Q}: both $P$ and $Q$ are satisfied simultaneously.
  \item \texttt{withFuel n P}: $P$ is satisfied when the computation is
        given at least $n$ units of fuel. (This constructor is used for
        fuel-parametric specifications and is not used in the main theorem
        statements.)
  \item \texttt{precond h P}: if the proposition $h$ holds, then $P$ is
        satisfied. This models preconditioned specifications.
\end{itemize}

\subsection{The Satisfaction Relation}

The satisfaction relation $m \satisfies P$ is defined by structural recursion
on $P$:

\begin{lstlisting}[caption={Satisfaction relation.},label={lst:sat}]
def satisfies (m : RustM s a) (init : s) :
    ProgramSpec s a -> Prop
  | .pureOutput P => exists  v, m init = .ok (v, init) /\ P v
  | .nocrash      => forall  e, m init /= .error e
  | .both P Q     => satisfies m init P /\ satisfies m init Q
  | .withFuel n P => satisfies m init P  -- simplified form
  | .precond h P  => h -> satisfies m init P

-- Notation: m at s |= P  means satisfies m s P
notation:50 m " at " s " |= " P => satisfies m s P
\end{lstlisting}

The \texttt{pureOutput} case deserves attention. The condition
\texttt{m init = .ok (v, init)} (not \texttt{.ok (v, s')} for some
\emph{different} $s'$) requires the state to be unchanged. This is always
satisfied by top-level function calls (where the state is \texttt{Unit}),
but it would fail for computations that write to a shared heap.
Since we do not model a shared heap, this simplification is sound.

\subsection{The Combinator Library}

The \texttt{Spec/Satisfies.lean} module proves 17 combinators. We list the
most frequently used ones:

\begin{lstlisting}[caption={Key ProgramSpec combinators (selection).},label={lst:combinators}]
-- Pure return: pure a satisfies pureOutput (. = a)
theorem sat_pure (a : a) (s : s) :
    pure a at s |= .pureOutput (. = a) := by
  simp [satisfies, pure_ok]

-- Sequential composition: if f returns a, and g a satisfies P, then f >>= g satisfies P
theorem sat_bind_pureOutput {m : RustM s a} {f : a -> RustM s b}
    (hm : m at s |= .pureOutput (. = a))
    (hf : f a at s |= P) :
    (m >>= f) at s |= P := by
  obtain (v, hv, rfl) := hm
  simp [bind_ok hm ...]

-- Conjunction introduction
theorem sat_both_intro
    (hP : m at s |= P) (hQ : m at s |= Q) :
    m at s |= .both P Q := (hP, hQ)

-- Monotonicity: if (. = v) implies P, then pureOutput (. = v) implies pureOutput P
theorem sat_pureOutput_mono {m : RustM s a}
    (h  : m at s |= .pureOutput (. = v))
    (hP : P v) :
    m at s |= .pureOutput P := by
  obtain (w, hw, rfl) := h; exact (w, hw, hP)
\end{lstlisting}

\texttt{sat\_pureOutput\_mono} is particularly important for the TypeScript
theorems (\S\ref{sec:typescript}): when the base-case execution equation is
blocked by the \texttt{Int.decLe @[extern]} issue, we can still derive
higher-level specifications from lower-level ones via this monotonicity lemma,
avoiding the need for \sorry in the derived theorems.

\begin{table}[t]
\centering
\caption{All 17 ProgramSpec combinators.}
\label{tab:combinators}
\small
\begin{tabular}{@{}ll@{}}
\toprule
\textbf{Combinator} & \textbf{Informal statement} \\
\midrule
\texttt{sat\_pure}              & \texttt{pure a} satisfies \texttt{pureOutput (. = a)} \\
\texttt{sat\_bind\_pureOutput}  & Compose: if $m$ returns $a$, then $m >>= f$ satisfies $f$'s spec \\
\texttt{sat\_both\_intro}       & Introduce conjunction \\
\texttt{sat\_both\_elim\_left}  & Eliminate left conjunct \\
\texttt{sat\_both\_elim\_right} & Eliminate right conjunct \\
\texttt{sat\_pureOutput\_mono}  & Weaken: \texttt{(.=v)} $\Rightarrow$ $P\,v$ yields $P$ \\
\texttt{sat\_nocrash\_of\_ok}   & If $m$ returns \texttt{ok}, it does not crash \\
\texttt{sat\_precond\_intro}    & Prove preconditioned spec by providing the precondition \\
\texttt{sat\_precond\_elim}     & Apply preconditioned spec \\
\texttt{sat\_rpanic\_nocrash}   & $\neg$: \texttt{rpanic} violates \texttt{nocrash} \\
\texttt{sat\_pureOutput\_eq}    & \texttt{pureOutput (. = v)} is extensionally equal \\
\texttt{sat\_map\_pureOutput}   & Functor: map preserves \texttt{pureOutput} \\
\texttt{sat\_bind\_nocrash}     & Bind: if $m$ and $f$ both don't crash, neither does $m >>= f$ \\
\texttt{sat\_precond\_True}     & If precondition is \texttt{True}, it can be discharged \\
\texttt{sat\_both\_comm}        & Commutativity of \texttt{both} \\
\texttt{sat\_withFuel\_intro}   & Introduce fuel-parametric spec \\
\texttt{sat\_withFuel\_elim}    & Eliminate fuel-parametric spec \\
\bottomrule
\end{tabular}
\end{table}

\subsection{RustM-Level Examples}

Before connecting \ProgramSpec to the LLBC embedding, it is useful to see it
applied to hand-written monadic terms. The \texttt{Spec/Examples.lean} module
contains eight such examples (E1--E8) to validate the combinator library in
isolation from the interpreter.

\begin{lstlisting}[caption={E1: rustAdd satisfies its specification.},label={lst:e1}]
-- Direct monadic definition of addition with overflow check
def rustAdd (x y : Int) : RustM Unit Int :=
  if minI32 <= x + y /\ x + y <= maxI32 then pure (x + y)
  else rpanic .arithmeticOverflow

theorem e1_add_correct (x y : Int)
    (h : minI32 <= x + y /\ x + y <= maxI32) :
    rustAdd x y at () |= .pureOutput (. = x + y) := by
  simp [rustAdd, satisfies, if_pos h]
\end{lstlisting}

\begin{lstlisting}[caption={E5: rustMax with two-case specification.},label={lst:e5}]
def rustMax (a b : Int) : RustM Unit Int :=
  pure (if a >= b then a else b)

theorem e5_max_ge (a b : Int) (h : a >= b) :
    rustMax a b at () |= .pureOutput (. = a) := by
  simp [rustMax, satisfies, if_pos h]

theorem e5_max_lt (a b : Int) (h : ¬(a >= b)) :
    rustMax a b at () |= .pureOutput (. = b) := by
  simp [rustMax, satisfies, if_neg h]

theorem e5_max_both (a b : Int) :
    rustMax a b at () |=
      .both (.pureOutput (. = if a >= b then a else b))
            .nocrash := by
  refine sat_both_intro ?_ ?_
  . simp [rustMax, satisfies]
  . simp [rustMax, satisfies]
\end{lstlisting}

These examples exercise the full combinator library and establish confidence
that the specification language is adequate before the LLBC interpreter is
introduced.

\subsection{Language Independence}

The \ProgramSpec type contains no reference to \texttt{LLBCFunDef},
\texttt{LLBCStmt}, \texttt{evalStmtFuel}, \texttt{TSFunDef}, or any other
language-specific construct. It is parametric in $\sigma$ (the state type)
and $\alpha$ (the value type). The satisfaction relation applies to any
\texttt{RustM}~$\sigma$~$\alpha$ computation.

This parametricity is the mechanism by which cross-language unification
theorems become possible. When the Rust evaluator \texttt{evalFun} and the
TypeScript evaluator \texttt{evalTSFun} both return computations of type
\texttt{RustM Unit Value}, the same \ProgramSpec specification
\texttt{.pureOutput (. = .int (if a $\ge$ b then a else b))} can be stated
and proved for both, and the two proofs can be composed into a unification
theorem without any additional logical infrastructure.
